% FDUletter_example.tex - an example latex file to illustrate FDUletter.cls
% OUCletter_example.tex - an example latex file to illustrate OUCletter.cls
%
% Template by Brian Wood (brian.wood@oregonstate.edu).  Please feel free to send suggestions for changes; this template/cls is not exactly elegantly done!
% Modified by Huang Weiran (huangweiran1998@outlook.com) to fit the need of FDU students.
% Modified by Tim (qt@stu.ouc.edu.cn) to fit the need of OUC students.


\documentclass[12pt]{OUCletter}
\usepackage{fontspec} 
\usepackage{tikz} 
\usepackage{xcolor}
\usepackage{hyperref}

\hypersetup{
  colorlinks=true,        % 让链接用颜色而不是方框标出
  linkcolor=black,        % 内部链接颜色（例如目录、交叉引用）
  citecolor=blue,         % 文献引用颜色
  urlcolor=blue, % 超链接颜色，保持和 scured 一致
  filecolor=magenta,      % 文件链接颜色
  pdfborder={0 0 0},      % 移除链接的边框
  breaklinks=true,        % 允许长链接换行
  bookmarksopen=true,     % 打开 PDF 书签
  bookmarksnumbered=true  % 书签带编号
}

	\definecolor{scured}{RGB}{1, 45, 114}
\usepackage{lipsum}
\usepackage{fancyhdr}
\usepackage{lastpage}
\usepackage{eso-pic}
%
% This section is just a bunch of busywork so that the second and following pages read ``Page X of Y''
\pagestyle{fancy}
\fancyhf{}
\renewcommand{\headrulewidth}{0pt}
\renewcommand{\footrulewidth}{0pt}
\rhead{Page \thepage \hspace{1pt} of \pageref{LastPage}}

\newcommand{\universityname}{University of Arizona}
\newcommand{\departmentname}{Department of Computer Science}
%
%
% Set custom font here. Comment this line out if you do not have a Cambria font (originally included with this template) installed; computer modern (or whatever your current default font is) will be substituted.
%
%\setmainfont{[Cambria.ttf]}[BoldFont  = [CambriaBold.ttf], ItalicFont  = [CambriaItalic.ttf], BoldItalicFont = [CambriaBoldItalic.ttf] ]

% \newcommand{\watermark}[3]{\AddToShipoutPictureBG{
% \parbox[b][\paperheight]{\paperwidth}{
% \vfill%
% \centering%
% \begin{tikzpicture}
%     \path (0,0) -- (\paperwidth,\paperheight);
%     \node[opacity=.05] at (current page.center)
%     {\includegraphics[width=0.55\textwidth]{background-OUC.pdf}};   % sculogo watermark
%     \end{tikzpicture}
% \vfill}}}

% The material below is a whole big dang thing whose purpose is just to set up a fixed coordinate system for \tikz so that you can put the Department or School address in the upper right-hand side without it moving all around every time you change something in the page.  I think it works.
% Defining a new coordinate system for the page:
%
% --------------------------
% |(-1,1)    (0,1)    (1,1)|
% |                        |
% |(-1,0)    (0,0)    (1,0)|
% |                        |
% |(-1,-1)   (0,-1)  (1,-1)|
% --------------------------
\makeatletter
\def\parsecomma#1,#2\endparsecomma{\def\page@x{#1}\def\page@y{#2}}
\tikzdeclarecoordinatesystem{page}{
	\parsecomma#1\endparsecomma
	\pgfpointanchor{current page}{north east}
	% Save the upper right corner
	\pgf@xc=\pgf@x%
	\pgf@yc=\pgf@y%
	% save the lower left corner
	\pgfpointanchor{current page}{south west}
	\pgf@xb=\pgf@x%
	\pgf@yb=\pgf@y%
	% Transform to the correct placement
	\pgfmathparse{(\pgf@xc-\pgf@xb)/2.*\page@x+(\pgf@xc+\pgf@xb)/2.}
	\expandafter\pgf@x\expandafter=\pgfmathresult pt
	\pgfmathparse{(\pgf@yc-\pgf@yb)/2.*\page@y+(\pgf@yc+\pgf@yb)/2.}
	\expandafter\pgf@y\expandafter=\pgfmathresult pt
}
\makeatother


%%------------------------------------------------------------------------------%%
%

%%%%%%%%%%% Put Recommender(your mentor) Information Here %%%%%%%%%%%
%
\def\name{\textbf{Jie Gao} \\[10pt]
	Malone Postdoctoral Fellow \\
	Department of Computer Science \\
    Malone Center for Engineering in Healthcare \\
    Center for Language and Speech Processing \\
    Johns Hopkins University \\
    jgao77@jh.edu \\
    \href{}{https://gaojie058.github.io/} \\
    
}
%
% List your degree(s), licences, etc. here if you like.
%\def\What{, Your degrees, etc.} 
%
% Set the name of your Department or School here
% I honestly don't know why the negative spacing is necessary to get the alignment of the first line correct.  This must be a ``\tikz'' thing.
%%%%%%%%%%%%%%%%%%  School or Department %%%%%%%%%%%%%%%
\def\Where{\hspace{-1.2mm}\textbf{\color{scured}%{scured}
		{Department of Computer Science}, \\
        {Johns Hopkins University} 
}} 

%%%%%%%%%%%%  Additional Contact Information %%%%%%%%%%%
%
% Set your preferred primary contact address here.
\def\Address{Malone Hall, 3400 N Charles St} 
%
\def\CityZip{Baltimore, MD 21218, USA %\\
%	P.R.China
} 
%
% Set your OUC e-mail here
\def\Email{\textbf{\color{scured}E-mail}: jgao77@jh.edu}
%
% Set your preferred contact number here
\def\TEL{\textbf{\color{scured}Phone}: +1 (410)710-0555}
%
% Set your department or personal website here
\def\URL{\textbf{\color{scured}Website}: {https://gaojie058.github.io/}}
%

%%%%%%%%%%%%%%%%  Footer information  %%%%%%%%%%%%%%%%%%
%
%  The next line is for your college, used as a footer.  If you prefer not to have this, just comment out these lines in favor of the line labeled "[[Alternate]]" below
\def\school{\small{
		OUC $\cdot$
		~Department of xxxxx $\cdot$
		~No. 238 Songling Road $\cdot$
		~Qingdao, China P.R} } 
% \def\school{~}  % [[Alternate]]
%

%%%%%%%%%%%%%%%%%%%%%  Signature line  %%%%%%%%%%%%%%%%%%%%%
%
% Set your signature line here.
% One can add a signature image in a PDF file using the following code; this requires a file called "signature_block.pdf" to be installed in the same folder as the .tex file.  The vertical spacing (\vspace) and the scaling will have to be adjusted to get things to look correct for your particular signature image. Alternatively, comment out the following line in favor of the one labeled "[[Alternate]]" if you want to sign a paper copy of the letter.
%
\signature{
	\vspace{2mm}\includegraphics[scale=0.8]{signature.pdf}\\\vspace{2mm}
	\name
}
%\signature{\name}  % [[Alternate]]
%%------------------------------------------------------------------------------%%



% This block sets up the address on the right-hand side of the header. 
%
% The following lines just compile the information you set up into the LaTex letter variable "address" for later use.
%
%The following command "clears out" the default address so that it can be better set using \tikz
\address{}

\def\newaddress{
	\Where\\ 
	\Address\\ 
	\CityZip\\ 
	\TEL\\ 
	\Email\\ 
	\URL 
}
%
%%%%%%%%%%%  DATE  %%%%%%%%%%%%%%%%%%%%%%%%%
% If you want a date different from the current date, comment out the next line in favor of the line labeled "[[Alternate]]".  The ``\vspace{10mm}'' just moves the date down a tiny bit so it doesn't interfere with the header.  This can be adjusted to your preference.
%
\date{\vspace{10mm} \today}
%\date{\vspace{10mm} 20 September 2020}  %[[Alternate]]
%
%%%%%%%%%%% Set the subject here if there is one  %%%%
%\subject{Stuff} % optional subject line



%%-----------------------------------------------------------------------------------------%%
%
\begin{document}
	%
	%
	%%%%%%%%  The "To" address goes here.
	%
	\begin{letter}{
			\textit{Faculty Search Committee} \\
            \departmentname \\
            \universityname
		%	IOP Publishing%\\ 
			%SomeTown, SomeState 					       				  		 ~~SomeZip
		}

		\begin{tikzpicture}[remember picture,overlay,,every node/.style={anchor=center}]
		\node[text width=7cm] at (page cs:0.5,0.73){\small \newaddress};
		\end{tikzpicture} 
		
		%%%%%%  The ``opening'' is just the method of address you would like to use at the start of the letter.
		%
		\opening{Dear Members of Search Committee,}
		
		
		%%%%%%%%%% Body of letter   %%%%%%%%%%%%%%
		% Remove it if you do not want watermark
		% \watermark{}{}{}
		
		%%%%%%%%%%% 	Text   %%%%%%%%%%%%%%
        It is with great enthusiasm that I apply for the \textbf{Tenure-Track Assistant Professor} position in \href{https://arizona.csod.com/ux/ats/careersite/4/home/requisition/24254?c=arizona&sq=req24254}{``AI and Data Science"} at \textbf{\universityname}. I am currently a \href{https://malonecenter.jhu.edu/people-2/members/}{Malone Postdoctoral Fellow} at Johns Hopkins University, where my work focuses on \textbf{human-AI collaboration}, \textbf{computational social science} and \textbf{responsible and trustworthy AI}.  
        I believe I am a strong candidate for this position due to my interdisciplinary background, shaped by both my education and research training in \textbf{Human-Computer Interaction (HCI)} and \textbf{Natural Language Processing (NLP)}, and because of my ability to generalize research problems across domains in ways that can contribute meaningfully to \universityname’s research mission.

        First, my research on analyzing unstructured data (e.g., conversations, interviews) has wide applicability across many domains and addresses important societal problems, including those in healthcare, public policy, and psychology research. For example, in clinical settings, conversations between doctors and patients contain rich information that can improve diagnoses and raise the quality of care. Yet, discovering insights from unstructured data remains difficult because of its scale, complexity, nuance, and lack of structure. This leads to a central question in my work: \textit{How can we design AI systems that strengthen human capabilities in analyzing complex unstructured data?} To address this, my projects \href{https://dl.acm.org/doi/abs/10.1145/3617362}{CoAIcoder} (TOCHI 2023), \href{https://dl.acm.org/doi/abs/10.1145/3613904.3642002}{CollabCoder} (CHI 2024), and \href{https://arxiv.org/abs/2501.00775}{MindCoder} (2025) develop AI systems to support efficient, transparent, and trustworthy data analysis. To support real world use, two of them are openly accessible through open source or public deployment \href{https://gaojie058.github.io/CollabCoder/}{CollabCoder} and \href{https://mindcoder.ai/}{MindCoder}. As some of the earliest investigations in this area, my published works at TOCHI and CHI have been well cited and impactful, which have inspired many follow up studies.


         Second, I believe I am a strong fit for this position because I can expand the boundaries of my research across disciplines, which enables me to collaborate with researchers in other domains and to create novel solutions for complex technical and societal challenges. For example, one of my recent interdisciplinary projects at the intersection of HCI and Software Engineering (SE), \href{https://arxiv.org/abs/2504.04553}{``Understanding Codebase like a Professional!"}, explores how insights from programmers can inform the design of AI-powered tools for codebase comprehension. In another example, my mentored student's \href{https://aclanthology.org/2025.hcinlp-1.20.pdf}{position paper} generalize social science research methods for application in subjective data annotation in NLP domain. These demonstrate that my research can be generalized across disciplines in both its methodology and its mindset.
        
        Lastly, I have also established professional connections with researchers from institutions such as MIT, Singapore–MIT Alliance for Research and Technology (SMART), University of Notre Dame, National University of Singapore (NUS), Singapore Management University (SMU), The Pennsylvania State University, and North Carolina State University. My previous role at SMART, which also gave me first hand experience collaborating with industry partners such as Singapore Changi Airport and Synergy Marine Group, further equipped me with the ability to translate research into industry relevant problems.
        
        Thank you for considering my application. I am confident that my research and interdisciplinary collaborations can contribute meaningfully to the \universityname’s mission, and I would also be honored to discuss how my research and teaching align with the \departmentname’s research vision.

        %  

		\closing{Sincerely yours,}	
	\end{letter}
	
\end{document}
%
%%-----------------------------------------------------------------------------------------%%